\bigskip

% course 6, page 1
% \section*{Curs 6 - Pagina 1}


\begin{center}
\fbox{\textbf{CURS 6}}\\[2pt]
\underline{(ALGEBRA)}
\end{center}
\begin{flushright}
\textit{8-XI-94}
\end{flushright}

\bigskip
\noindent\underline{\textbf{Ordinul unui element}}

\bigskip
\noindent\underline{Anexa I}: -- Elemente de aritmetică în $\mathbb{N}$.

\bigskip
\noindent\underline{Sistem Peano}: $(A,a,s)$ ; $a\in A$; $s:A\to A$,\footnote{în original apare ,,$s:A\to B$''; corectat aici la $A\to A$, cum cere definiția unui sistem Peano.} $A\neq\varnothing$

\begin{enumerate}
\item[1)] $s(x)\neq a$ $(\forall)\,x\in A$, \ $\big(s(A)=A\setminus\{a\}\big)$.
\item[2)] $s$ injectivă.
\item[3)] Dacă $M\subseteq A$ are pp.\ $\begin{cases} a\in M \\ x\in M \Rightarrow s(x)\in M \end{cases} \Rightarrow M=A$.
\end{enumerate}

\bigskip
\noindent\fbox{T1}. \ Dacă $(A,a,s)$, $(B,b,t)$ sisteme Peano $\Rightarrow$

\noindent $\exists!\, f:A\to B$ a.î.\ $f(a)=b$, $f\circ s = t\circ f$.

\noindent Mai mult,\footnote{fraza exactă e neclară în original.} $f$ e bijectivă.

\[
\begin{CD}
A @>s>> A\\
@VfVV @VVfV\\
B @>t>> B
\end{CD}
\]

\bigskip
\noindent Fie $(\mathbb{N},0,s)$; \ $s:\mathbb{N}\to\mathbb{N}$.

\noindent funcţ.\ succesor.

\noindent $s(x)=x'$ -- succesorul lui $x$.

\bigskip
\noindent N.\ Becheanu -- Alg.\ ptr.\ perf.\ profesorale\footnote{referință bibliografică; titlul exact e reconstituit din formă compactată.}

\bigskip
\noindent\fbox{T2}: $\exists!\, e:\mathbb{N}\times\mathbb{N}\to\mathbb{N}$ \ $e\overset{\text{not}}{=}+$

\noindent $A_1$: $e(m,0)=m$ $(\forall)\,m\in\mathbb{N}$.

\noindent $A_2$: $e(m,n')=e(m,n)'$ $(\forall)\,m,n\in\mathbb{N}$ \quad $m+n'=(m+n)'$.

\bigskip
\noindent Rezultă $m+p=m+n \Rightarrow p=n$.

\[
(\mathbb{N},+,0) = \text{monoid comutativ cu simplificare.}
\]

\bigskip
\noindent\fbox{T}: $\exists!\, \psi:\mathbb{N}\times\mathbb{N}\to\mathbb{N}$ a.î.

\noindent $\psi(m,0)=0$\footnote{în original pare ,,$=m$''; corectat la $0$, cum cere definiția standard a înmulțirii (notația alăturată ,,$m\cdot 0$'' e la fel $0$, nu $m$).} $(\forall)\,m\in\mathbb{N}$: \quad $m\cdot 0=0$

\noindent $\psi(m,n')=\psi(m,n)+m$ $(\forall)\,m,n\in\mathbb{N}$ \quad $mn'=mn+m$.

\bigskip
\noindent $(\mathbb{N},\cdot,1)$ = monoid comutativ cu simplificare.

\noindent $mn=mp \Rightarrow n=p$.

\noindent $m(n+p)=mn+mp$.

\bigskip
\noindent $(\mathbb{N},+,\cdot)$ = structură algebrică.

\noindent $(\mathbb{N},\leq)$ = structură de ordine.



% course 6, page 2
\section*{Curs 6 - Pagina 2}

\noindent $m<n$ dacă $\exists\, k\neq 0$ a.î.\ $m+k=n$.

\noindent $m\leq n \Longleftrightarrow m<n$ sau $m=n$.

\noindent reflexivă, antisimetrică, tranzitivă.

\bigskip
\noindent $(\forall)\,m,n\in\mathbb{N}$, una şi numai una din prop.\ următoare este adevărată:

\begin{enumerate}
\item[(1)] $m<n$
\item[(2)] $m=n$
\item[(3)] $m>n$.
\end{enumerate}

\bigskip
\noindent\fbox{T} \ $(\mathbb{N},\leq)$ mulţ.\ bine ordonată.

\noindent $(\forall)\,P\subseteq\mathbb{N}$, $P\neq\varnothing$, $\exists!\,a\in P$ a.î.\ $a\leq x$ $(\forall)\,x\in P$.

\[
(\mathbb{N},+,\cdot,\leq)
\]

\bigskip
\noindent\fbox{T} \ ,,$\leq$'' este compatibilă cu str.\ algebrică.

\[
m\leq n \Rightarrow
\begin{cases}
m+p\leq n+p, & (\forall)\,p\in\mathbb{N}.\\
mp\leq np, & (\forall)\,p\in\mathbb{N}^*.
\end{cases}
\]

\bigskip
\noindent\fbox{T} \ $(\forall)\,a,b\in\mathbb{N}$, $b\neq 0$, $\exists!\,q,r\in\mathbb{N}$ a.î.\ $a=bq+r$, $r<b$.

\noindent (împărţirea euclidiană (cu rest)).

\bigskip
\noindent $(\mathbb{N},\mid)$ ; \ ,,$\mid$''-rel.\ de divizibilitate. \hfill $1\leq a\leq b$

\noindent Dacă $a,b\in\mathbb{N}$, $a\mid b \Leftrightarrow \exists\,q\in\mathbb{N}$ a.î.\ $b=aq$.

\bigskip
\noindent $p>1$ s.n.\ \underline{prim} dacă $p\mid ab \Rightarrow p\mid a$ sau $p\mid b$.

\noindent (idem) \underline{ireductibil} (indescompozabil) dacă $d\mid p \Rightarrow d=1$ sau $d=p$.

\bigskip
\noindent\fbox{T} \ Prim $\Leftrightarrow$ ireductibil

\bigskip
\noindent\underline{Dem}: \ ,,$\Rightarrow$''

\noindent Pp.\ $p$=prim şi $d\mid p \Rightarrow \exists\,q\in\mathbb{N}$ a.î.\ $p=dq$.

\[
\Rightarrow p\mid dq \Rightarrow p\mid d \text{ sau } p\mid q \Rightarrow \underline{d=p\cdot q_1}
\]
\[
\Rightarrow qq_1=1 \Rightarrow d=1 \text{ sau } d=p.
\]

\bigskip
\noindent\underline{Def}: Dacă $a,b\in\mathbb{N}$, un nr.\ natural $d$ s.n.\ \underline{c.m.m.d.c.\ al}

\noindent lui $a,b$ dacă:

\begin{enumerate}
\item[(1)] $d\mid a$ şi $d\mid b$.
\item[(2)] $c\mid a$ şi $c\mid b \Rightarrow c\mid d$.
\end{enumerate}

